\subsection{The cofinite theory of finite groups: Problem 18.18}
\label{cand:18.18}
\problemstatement{18.18}

Let \(T_{\rm fin}\) consist of group-language sentences true
in every finite group, and \(T_{\rm cof}\) those true in
all but finitely many finite groups up to isomorphism.
The latter is equivalent to truth in all sufficiently
large finite groups. Both \(T_{\rm cof}\) and its
complement fail to be computably enumerable.
The input theorem is Mal'cev's undecidability of
\(T_{\rm fin}\) \cite{malcev-finite}: finite
countermodels can be enumerated, so \(T_{\rm fin}\)
is co-computably-enumerable and cannot also be
computably enumerable.

For a sentence \(\phi\), effectively relativize every
quantifier to the definable subgroup
\[
 C(x;a,b):\quad xa=ax\ \wedge\ xb=bx,
\]
renaming bound variables first. Write the resulting
formula as \(\phi^C(a,b)\). Formula induction gives
its truth exactly when \(C_G(a,b)\models\phi\).
Then \(P_\phi=\forall a\,\forall b\,\phi^C(a,b)\)
satisfies
\[
 \phi\in T_{\rm fin}\quad\Longleftrightarrow\quad
 P_\phi\in T_{\rm cof}.
\]
The forward direction is immediate. For the reverse,
if finite \(H\not\models\phi\), use \(H\times S_n\),
\(n\geq3\), and two generators of the symmetric factor.
Their common centralizer is \(H\times1\), giving
countermodels of unbounded order. This proves the
first non-enumerability assertion.

For the complement we need a fixed first-order
sentence \(\sigma\) whose finite models are exactly
the \(S_n\), \(n\geq5\). The following explicit
interpretation supplies it. Choose an involution
\(t\ne1\), let \(C_t\) be its conjugacy class,
and let \(Q(\boldsymbol u)\) say that four entries
of \(\boldsymbol u\) are distinct members of \(C_t\)
and pairwise noncommuting. Put
\[
 \begin{split}
 M(c,\boldsymbol u)&:\quad c\in C_t\ \wedge\
   \bigwedge_{i=1}^4(c=u_i\ \vee\ cu_i\ne u_ic),\\
 E(\boldsymbol u,\boldsymbol v)&:\quad
   \forall c\in C_t\
   \bigl(M(c,\boldsymbol u)\leftrightarrow
                    M(c,\boldsymbol v)\bigr).
 \end{split}
\]
On \(Q\)-tuples, \(E\) is equality of the incidence
sets and hence an equivalence relation. Call these
classes vertices. This is abbreviation only:
vertex quantifiers use four group variables and
vertex equality uses \(E\).
Require that there be at least five vertices; that
each \(c\in C_t\) be incident with exactly two
vertices; that every pair of distinct vertices have
exactly one common incident \(c\); that
\[
 E(c\boldsymbol u c^{-1},\boldsymbol u)
       \quad\Longleftrightarrow\quad
       \neg M(c,\boldsymbol u);
\]
and that the common centralizer of \(C_t\) be trivial.
These are finitely many fixed first-order assertions;
their conjunction, existentially quantifying \(t\),
defines \(\sigma\).

In a finite model, conjugation acts faithfully on
the finite vertex set: an element fixing all vertices
fixes their unique edges and therefore all \(C_t\).
Each \(c\) preserves its incident pair, moves both
endpoints by the displayed assertion, and fixes
every other vertex. Its action is that pair's
transposition. Every pair occurs, so the faithful
image is the full symmetric group.
Conversely, in \(S_n\), take \(t\) a transposition.
Four distinct pairwise intersecting two-element
subsets have a common point: the only alternative
for three is a triangle, which admits no fourth
distinct edge meeting all three. The sets defined
by \(M\) are therefore exactly the stars at points.
Every point supplies a four-tuple when \(n\geq5\).
All the assertions hold, including trivial
centralizer because transpositions generate
centerless \(S_n\). This proves the claim about
\(\sigma\).

We also need every finite \(H\) as a common
centralizer in \emph{every sufficiently large}
\(S_n\). Write \(h=|H|\), enumerate
\(H=\{h_0=1,h_1,\ldots,h_{h-1}\}\), and on
\(\Omega=H\times\{0,\ldots,h-1\}\) define
\[
 a(g,j)=(g,j+1\bmod h),\qquad b(g,j)=(gh_j,j).
\]
The fixed points of \(b\) form precisely level zero.
A commuting permutation must therefore have form
\((g,j)\mapsto(u(g),j)\) by commutation with \(a\);
commutation with \(b\) gives
\(u(gh_j)=u(g)h_j\), hence \(u(g)=u(1)g\).
The common centralizer is exactly the left regular
copy of \(H\). The action of \(\langle a,b\rangle\)
is transitive, by moving between levels to realize
all right multiplications.

For \(n\geq2h^2+3\), add a disjoint orbit of
\(k=n-h^2>h^2\) points, extending \(a,b\) by
a \(k\)-cycle and a transposition of consecutive
points. Their action there is \(S_k\), with
trivial centralizer. The two orbit sizes differ,
so no commuting permutation exchanges them.
The total common centralizer in \(S_n\) is
therefore \(H\), as required.

Finally form the effective sentence
\[
 R_\phi=(\neg\sigma)\ \vee\
             \exists a\,\exists b\,\neg\phi^C(a,b).
\]
If \(\phi\in T_{\rm fin}\), the second disjunct
is false in every finite group, and \(R_\phi\)
fails in every \(S_n\), \(n\geq5\). If a finite
\(H\) refutes \(\phi\), all nonsymmetric groups
satisfy the first disjunct, and every sufficiently
large symmetric group satisfies the second by
the construction. Thus
\[
 \phi\in T_{\rm fin}\quad\Longleftrightarrow\quad
 R_\phi\notin T_{\rm cof}.
\]
Computable enumerability of the complement would
give it for \(T_{\rm fin}\), a contradiction.
The syntactic transformation does not search for
\(H\); the witness is needed only for its correctness
proof. No asymptotic-density interpretation of
\emph{almost all} is asserted. Source:
\artifact{research/18.18-proof.md}.

